\section{Walkers}

\subsection{Sydonian Dragoons}

Sydonian Dragoons cross the battlefield like living engravings, bursting from clouds of sacred incense to plunge their crackling lances into the heart of a panicked enemy. They leave behind them only a wake of corpses stretching beyond the horizon. Mounted on the scientific masterpiece that is the Ironstrider walker, Sydonian Dragoons know neither doubt nor rest. The first Sydonians settled in a crater-strewn region of Mars bathed in an acid mist. If the Servitor grafted to each Ironstrider walker is virtually brainless, that is not the case of its rider, who is a seasoned Skitarii.

The Sydonian Dragoon has the \textbf{Lance (5+ / AP $-1$)} and \textbf{Infiltration} abilities.

\begin{unitstatstable}
20 cm & 5+ & +3 & 6 & -- & -- & -- & -- & --
\end{unitstatstable}

\unitimage{sydonian-dragoons.png}

\unitdivider

\subsection{Ironstrider Ballistarii}

The Ironstrider Ballistarii mount the same walkers as the Sydonian Dragoons. These machines are capable of crossing the most hostile terrain for years without pause. The single-task Servitor that operates the Ironstrider is often the first component to fail, its battered body finally yielding its last breath. On such occasions, the Ironstrider will continue its crossing until its last recorded destination. It can be repaired only by a Tech-Priest on an anti-grav stretcher, who will excise the dead flesh and replace it while the walker continues its progress. The Ironstrider will then resume its task as if nothing had happened. Thanks to the Heavy Autocannons at its fingertips and the enemy's weak points in its sights, a Ballistarius could fire through the hull of an enemy tank and hit the tank commander in the eye. Indeed, bringing down the enemy's principal vehicles is only a secondary objective compared with hunting its commanders, heretical priests, or xenos warlords, for silencing these leaders sows despair in the opposing army.

The Ironstrider Ballistarii have the \textbf{Fire on the Move} and \textbf{Infiltration} abilities.

\begin{unitstatstable}
20 cm & 5+ & +2 & 6 & Heavy Autocannon & 45 cm & 1 & 4+ & $-1$
\end{unitstatstable}

\unitimage{ironstrider-ballistarii.png}

\unitdivider

\subsection{Magos Dominus on Abeyant}

The Magos Dominus is a master of the Legio Cybernetica, of Electro-Priests, and of Servitors. A master in the art of war, he is able to read the battlefield with a talent worthy of a machine. They are often mounted on an Abeyant, a class of hovering vehicles in which the pilot's augmetic and life-support systems are directly connected by bionic links, so that the machine-vehicle becomes an extension of his own body.

He has the \textbf{HQ}, \textbf{Attached Character}, \textbf{Elite (1)}, \textbf{Regeneration (5+)}, \textbf{Mechanic}, and \textbf{Robotics Master} abilities. In addition his Grav-gun has the \textbf{Graviton} ability.

\begin{unitstatstable}
20 cm & 4+/5+f & +5 & Attached & Grav-gun & 45 cm & 1 & Graviton & $-2$
\end{unitstatstable}

\unitimage{magos-dominus-on-abeyant.png}

\unitdivider

\subsection{Arlatax}

Arlatax-class robots were a type of battle robot used by the Legio Cybernetica during the Great Crusade and the Horus Heresy. A variant of the ancient Conqueror class developed on Xana, its design influenced the Domitar class. Despite its excellent performance as a fast-moving shock unit, it was never widely accepted by the Mechanicum, for it was said that some of its components were not approved by Mars. That made repairs and overhauls difficult, and rumours also circulated that the machine spirit was prone to corruption. However, as the Heresy progressed, these robots were widely used.

It has the \textbf{Robot}, \textbf{Jump Packs}, and \textbf{Reflex Fire} abilities.

\begin{unitstatstable}
20 cm & 3+ & +5 & -- & Autocannon & 45 cm & 1 & 4+ & 0
\end{unitstatstable}

\unitimage{arlatax.png}

\unitdivider

\subsection{Castellan}

The Castellan is a versatile Battle-automata developed to be used in a wide range of situations on the battlefield. The two Power Fists in the standard configuration make it particularly effective in siege and anti-vehicle roles, and form an effective complement to the heavy armament.

It has the \textbf{Robot} and \textbf{Reflex Fire} abilities. Its weapon has the \textbf{Reduces Cover ($-3$)} ability.

\begin{unitstatstable}
15 cm & 4+ & +5 & -- & Flamer & 20 cm & 2 & 4+ & 0
\end{unitstatstable}

\unitimage{castellan.png}

\unitdivider

\subsection{Cataphract}

The Cataphract is a heavily armoured Battle-automata designed to be used in almost every combat situation. It possesses a whole range of weapons for various purposes, and although it is regarded by some commanders as suited to everything without having a precise speciality, its versatility in a wide range of roles makes it a highly favoured and popular machine.

It has the \textbf{Robot} and \textbf{Reflex Fire} abilities.

\begin{unitstatstable}
\SetCell[r=2]{c,m} 10 cm &
\SetCell[r=2]{c,m} 4+ &
\SetCell[r=2]{c,m} +2 &
\SetCell[r=2]{c,m} -- &
Lascannon &
60 cm &
1 &
5+ &
$-1$
\\
& & & &
Heavy Bolter &
45 cm &
1 &
5+ &
$-1$
\end{unitstatstable}

\unitimage{cataphract.png}

\unitdivider

\subsection{Colossus}

The Colossus is first and foremost a siege Battle-automata; the standard configuration combines a siege hammer to attack fortifications with lighter weapons to use against defending troops. It is widely used by Mechanicum forces as a less costly alternative to the Siege Robot. A few Colossus-class Battle-automata are maintained by the Ordo Reductor of the Adeptus Mechanicus, mainly to test new and experimental siege weapons.

The Colossus has the \textbf{Robot}, \textbf{Reflex Fire}, and \textbf{Damages Buildings (AP $-2$/1) in assault} abilities.

\begin{unitstatstable}
15 cm & 4+ & +4 & -- & Multi-melta & 20 cm & 1 & 3+ & $-2$
\end{unitstatstable}

\unitimage{colossus.png}

\unitdivider

\subsection{Conqueror}

The Conqueror Battle-automata is designed to be used on the battlefield in situations where strong armed resistance is expected. A major point of its design is the balance and combination of firepower with heavy armour. It is widely used in assaults and the elimination of vehicles and Dreadnoughts.

It has the \textbf{Robot} and \textbf{Reflex Fire} abilities.

\begin{unitstatstable}
\SetCell[r=2]{c,m} 10 cm &
\SetCell[r=2]{c,m} 4+ &
\SetCell[r=2]{c,m} +4 &
\SetCell[r=2]{c,m} -- &
Heavy Bolter &
45 cm &
1 &
5+ &
$-1$
\\
& & & &
Autocannon &
45 cm &
1 &
4+ &
0
\end{unitstatstable}

\unitimage{conqueror.png}

\unitdivider

\subsection{Domitar}

The last type of battle robot to enter production within the Mechanicum before the Horus Heresy, the Domitar was a variant of the sophisticated ancient Conqueror-class robot. More imposing than the Castellax class alongside which it would serve, the Domitar was also faster, designed to cross the battlefield at high speed and break through enemy positions. To that end it was given a heavily armoured, shock-resistant chassis and is equipped with grav-hammers. As a line-breaker, the Domitar was without equal in the Mechanicum's arsenal, and its tactical flexibility was further reinforced.

It has the \textbf{Damages Buildings (AP $-4$/1) in assault} and \textbf{Robot} abilities.

\begin{unitstatstable}
15 cm & 3+/5+f & +5 & -- & Grav-hammers & -- & -- & -- & --
\end{unitstatstable}

\unitimage{domitar.png}

\unitdivider

\subsection{Thanatar}

The Thanatar is a heavy Siege-automata, designed as a mobile artillery platform rather than as a general combat unit. It is armed with a Hellex-pattern plasma mortar, a terrifying weapon capable of launching dense spheres of flaming plasma on arcing trajectories over defensive barriers and into the heart of enemy fortifications, incinerating everything within its blast radius with a liquid tide of solar fire. The Thanatar's chassis, built to house the enormous weapon and the power systems and fuel reserves needed to feed it. Its structure is also heavily reinforced, both to withstand the firing stresses of its main weapon and to protect it against external attacks. Unlike many other Battle-automata produced to add the Mechanicum's power to that of the Great Crusade, the origins and lineage of the Thanatar pattern remain uncertain.

The Thanatar has the \textbf{Robot} ability. Its weapon has the \textbf{Reduces Cover ($-1$)} and \textbf{Artillery} abilities.

\begin{unitstatstable}
10 cm & 3+ & +2 & -- & Plasma Mortar & 60 cm & 1 & 3+ & $-1$
\end{unitstatstable}

\unitimage{thanatar.png}

\unitdivider

\subsection{Thanatar Calix}

The Calix-class Thanatar is said to have been developed by the Sollex Myrmidon and Omega-Shevar alliances of the Ordo Reductor alongside the Legio Cybernetica, and to combine in a single weapon the deadly technologies of those three factions. Calix-class Thanatars are intended for the vanguard of assault warfare, able to break the strongest fortifications and carry out targeted strikes to neutralize defensive strongpoints. Its principal weapon is the Sollex Heavy Lascannon, and these engines are so precious that, in the rare cases where they are deployed, many other lesser combat automata are deployed to protect them.

The Thanatar Calix has the \textbf{Robot} and \textbf{Reflex Fire} abilities. In addition its Grav Ram has the \textbf{Graviton} and \textbf{Damages Buildings (AP $-3$/1)} abilities.

\begin{unitstatstable}
\SetCell[r=2]{c,m} 10 cm &
\SetCell[r=2]{c,m} 3+ &
\SetCell[r=2]{c,m} +3 &
\SetCell[r=2]{c,m} -- &
Grav Ram &
20 cm &
1 &
Graviton &
$-2$
\\
& & & &
Lascannon &
60 cm &
1 &
5+ &
$-1$
\end{unitstatstable}

\unitimage{thanatar-calix.png}

\unitdivider

\subsection{Vorax}

The Vorax are hunter-killer units, Battle-automata created from the ancient and venerated Crusader pattern, which is said to date back to the Dark Age of Technology and can be found on all the great Forge Worlds of the Imperium. Created on pre-Imperial Mars to fulfil both the roles of weapon of destruction and exterminator of malevolent machines and mutant vermin, their Cybernetica cortex engrams are known to be particularly predatory and vicious examples. The bestial instincts that sleep in their bodies of steel and ceramite are such that many commands are built into their control programming to prevent these tireless hunting machines from becoming like the renegades they attack in the shadows of Mars's great forges. Sometimes they are also sent by their masters to cut down surplus populations in times of famine or plague, their hunts stopping and their guns falling silent abruptly. It is therefore no surprise that the Vorax are superstitiously feared by the indentured workers of the Mechanicus's distant domains and by neophyte adepts of the priesthood of the Omnissiah.

The Vorax has the \textbf{Robot} and \textbf{Reflex Fire} abilities.

\begin{unitstatstable}
15 cm & 4+ & +3 & -- & Rotor Cannon & 30 cm & 3 & 5+ & $-1$
\end{unitstatstable}
\unitimage{vorax.png}

